% GENERATED FILE. Edit canonical lessons, then run npm run generate:books.
% canonical-source-hash: fnv1a64:cf3e6b5541d645eb
% canonical-lessons: TA-C501-caccaravu, TA-C501-takararu, TA-C501-manastapam, TA-C501-pinanku, TA-C501-vivakarattu

\chapter{Family Quarrels}
\label{ch:ta-family-quarrels}
\hlchaptermodality{501}

\begin{chapteropening}
\textbf{By the end of this chapter:} \emph{I can talk about squabbles, disputes, ill feeling, sulking and divorce.}

Complete the last lesson of chapter 501: I can talk about squabbles, disputes, ill feeling, sulking and divorce.
\end{chapteropening}

\section[\texorpdfstring{\ta{சச்சரவு} (caccaravu)}{சச்சரவு (caccaravu)}]{\ta{சச்சரவு} (caccaravu) — a squabble}

\label{lesson:TA-C501-caccaravu}



\begin{quote}
Before the new one: say the Tamil for a separate household, then the Tamil for to get along.
\end{quote}

\subsection*{You'll want to know: \ta{சச்சரவு}}
\textbf{\ta{சச்சரவு}} — \emph{caccaravu} — \textquotedblleft{}a squabble\textquotedblright{}.

\begin{quote}
\ta{சிறிய} \ta{விஷயத்துக்கு} \ta{அண்ணன்} \ta{தம்பிக்குள்} \ta{சச்சரவு} \ta{வந்தது}.
\end{quote}

\begin{quote}
\emph{ciṟiya viṣayattukku aṇṇaṉ tampikkuḷ caccaravu vantatu.}
\end{quote}

\begin{quote}
The brothers squabbled over a small matter.
\end{quote}

\ta{சண்டை} \ta{சச்சரவு}, fights and squabbles, is a common pair.

\begin{cousinweb}
The origin of \textbf{\ta{சச்சரவு}} (\emph{caccaravu}) is not settled; we could not find it in DEDR.
\end{cousinweb}

\begin{grammarlens}[title={What you've built}]
One more word for describing what you lived and felt.
\end{grammarlens}

\subsection*{Your turn}
\begin{itemize}
\raggedright
  \item \emph{Say it:} \emph{caccaravu}
  \item \emph{Say it:} \emph{caccaravu}, once more
  \item \emph{Say it:} \emph{ciṟiya viṣayattukku aṇṇaṉ tampikkuḷ caccaravu vantatu.} — The brothers squabbled over a small matter.
  \item \emph{Recall:} say the Tamil for a separate household, then the Tamil for to get along, then say \emph{caccaravu} again
\end{itemize}

\begin{tcolorbox}[breakable,colback=teal!4,colframe=teal!35!black,title={Before you move on}]
What does \ta{சச்சரவு} mean? (\textquotedblleft{}A squabble\textquotedblright{}.) Say it once more.
\end{tcolorbox}
\section[\texorpdfstring{\ta{தகராறு} (takarāṟu)}{தகராறு (takarāṟu)}]{\ta{தகராறு} (takarāṟu) — a dispute; a row}

\label{lesson:TA-C501-takararu}



\begin{quote}
Before the new one: say the Tamil for to get along, then the Tamil for a squabble.
\end{quote}

\subsection*{You'll want to know: \ta{தகராறு}}
\textbf{\ta{தகராறு}} — \emph{takarāṟu} — \textquotedblleft{}a dispute; a row\textquotedblright{}.

\begin{quote}
\ta{நிலம்} \ta{பற்றி} \ta{இரண்டு} \ta{குடும்பங்களுக்கும்} \ta{தகராறு}.
\end{quote}

\begin{quote}
\emph{nilam paṟṟi iraṇṭu kuṭumpaṅkaḷukkum takarāṟu.}
\end{quote}

\begin{quote}
The two families have a dispute over land.
\end{quote}

\ta{தகராறு} \ta{செய்}, to pick a quarrel.

\begin{cousinweb}
\textbf{\ta{தகராறு}} (\emph{takarāṟu}) is usually traced to Arabic \emph{takrār}, \textquotedblleft{}repetition; altercation\textquotedblright{}, which reached Tamil through Hindustani; our sources do not document the path.
\end{cousinweb}

\begin{grammarlens}[title={What you've built}]
One more word for describing what you lived and felt.
\end{grammarlens}

\subsection*{Your turn}
\begin{itemize}
\raggedright
  \item \emph{Say it:} \emph{takarāṟu}
  \item \emph{Say it:} \emph{takarāṟu}, once more
  \item \emph{Say it:} \emph{nilam paṟṟi iraṇṭu kuṭumpaṅkaḷukkum takarāṟu.} — The two families have a dispute over land.
  \item \emph{Recall:} say the Tamil for to get along, then the Tamil for a squabble, then say \emph{takarāṟu} again
\end{itemize}

\begin{tcolorbox}[breakable,colback=teal!4,colframe=teal!35!black,title={Before you move on}]
What does \ta{தகராறு} mean? (\textquotedblleft{}A dispute; a row\textquotedblright{}.) Say it once more.
\end{tcolorbox}
\section[\texorpdfstring{\ta{மனஸ்தாபம்} (maṉastāpam)}{மனஸ்தாபம் (maṉastāpam)}]{\ta{மனஸ்தாபம்} (maṉastāpam) — ill feeling; estrangement}

\label{lesson:TA-C501-manastapam}



\begin{quote}
Before the new one: say the Tamil for a squabble, then the Tamil for a dispute.
\end{quote}

\subsection*{You'll want to know: \ta{மனஸ்தாபம்}}
\textbf{\ta{மனஸ்தாபம்}} — \emph{maṉastāpam} — \textquotedblleft{}ill feeling; estrangement\textquotedblright{}.

\begin{quote}
\ta{ஒரு} \ta{சிறிய} \ta{பேச்சால்} \ta{அவர்களுக்குள்} \ta{மனஸ்தாபம்} \ta{ஏற்பட்டது}.
\end{quote}

\begin{quote}
\emph{oru ciṟiya pēccāl avarkaḷukkuḷ maṉastāpam ēṟpaṭṭatu.}
\end{quote}

\begin{quote}
A small remark caused ill feeling between them.
\end{quote}

Written with the grantha \ta{ஸ}.

\begin{cousinweb}
\textbf{\ta{மனஸ்தாபம்}} (\emph{maṉastāpam}) is borrowed from Sanskrit \emph{manastāpa}, \textquotedblleft{}pain or anguish of the mind\textquotedblright{} (Monier-Williams); Tamil uses it for a rift between people.
\end{cousinweb}

\begin{grammarlens}[title={What you've built}]
One more word for describing what you lived and felt.
\end{grammarlens}

\subsection*{Your turn}
\begin{itemize}
\raggedright
  \item \emph{Say it:} \emph{maṉastāpam}
  \item \emph{Say it:} \emph{maṉastāpam}, once more
  \item \emph{Say it:} \emph{oru ciṟiya pēccāl avarkaḷukkuḷ maṉastāpam ēṟpaṭṭatu.} — A small remark caused ill feeling between them.
  \item \emph{Recall:} say the Tamil for a squabble, then the Tamil for a dispute, then say \emph{maṉastāpam} again
\end{itemize}

\begin{tcolorbox}[breakable,colback=teal!4,colframe=teal!35!black,title={Before you move on}]
What does \ta{மனஸ்தாபம்} mean? (\textquotedblleft{}Ill feeling; estrangement\textquotedblright{}.) Say it once more.
\end{tcolorbox}
\section[\texorpdfstring{\ta{பிணங்கு} (piṇaṅku)}{பிணங்கு (piṇaṅku)}]{\ta{பிணங்கு} (piṇaṅku) — to sulk; to fall out}

\label{lesson:TA-C501-pinanku}



\begin{quote}
Before the new one: say the Tamil for a dispute, then the Tamil for ill feeling.
\end{quote}

\subsection*{You'll want to know: \ta{பிணங்கு}}
\textbf{\ta{பிணங்கு}} — \emph{piṇaṅku} — \textquotedblleft{}to sulk; to fall out\textquotedblright{}.

\begin{quote}
\ta{சண்டைக்குப்} \ta{பிறகு} \ta{மனைவி} \ta{கணவரிடம்} \ta{பிணங்கினாள்}.
\end{quote}

\begin{quote}
\emph{caṇṭaikkup piṟaku maṉaivi kaṇavariṭam piṇaṅkiṉāḷ.}
\end{quote}

\begin{quote}
After the quarrel the wife sulked with her husband.
\end{quote}

Past (I …): \ta{பிணங்கினேன்}, \emph{piṇaṅkiṉēṉ}.

The one sulked at takes -\ta{இடம்}.

\begin{cousinweb}
\textbf{\ta{பிணங்கு}} (\emph{piṇaṅku}) is native Dravidian (DEDR 4160), in the entry of \ta{பிணை} \emph{piṇai}, \textquotedblleft{}to entwine, to unite\textquotedblright{}.
\end{cousinweb}

\begin{grammarlens}[title={What you've built}]
One more word for describing what you lived and felt.
\end{grammarlens}

\subsection*{Your turn}
\begin{itemize}
\raggedright
  \item \emph{Say it:} \emph{piṇaṅku}
  \item \emph{Say it:} \emph{piṇaṅku}, once more
  \item \emph{Say it:} \emph{caṇṭaikkup piṟaku maṉaivi kaṇavariṭam piṇaṅkiṉāḷ.} — After the quarrel the wife sulked with her husband.
  \item \emph{Recall:} say the Tamil for a dispute, then the Tamil for ill feeling, then say \emph{piṇaṅku} again
\end{itemize}

\begin{tcolorbox}[breakable,colback=teal!4,colframe=teal!35!black,title={Before you move on}]
What does \ta{பிணங்கு} mean? (\textquotedblleft{}To sulk; to fall out\textquotedblright{}.) Say it once more.
\end{tcolorbox}
\section[\texorpdfstring{\ta{விவாகரத்து} (vivākarattu)}{விவாகரத்து (vivākarattu)}]{\ta{விவாகரத்து} (vivākarattu) — a divorce}

\label{lesson:TA-C501-vivakarattu}



\begin{quote}
Before the new one: say the Tamil for ill feeling, then the Tamil for to sulk.
\end{quote}

\subsection*{You'll want to know: \ta{விவாகரத்து}}
\textbf{\ta{விவாகரத்து}} — \emph{vivākarattu} — \textquotedblleft{}a divorce\textquotedblright{}.

\begin{quote}
\ta{பல} \ta{வருடச்} \ta{சண்டைக்குப்} \ta{பிறகு} \ta{அவர்கள்} \ta{விவாகரத்து} \ta{செய்தார்கள்}.
\end{quote}

\begin{quote}
\emph{pala varuṭac caṇṭaikkup piṟaku avarkaḷ vivākarattu ceytārkaḷ.}
\end{quote}

\begin{quote}
After years of fighting they divorced.
\end{quote}

\ta{விவாகம்} \textquotedblleft{}marriage\textquotedblright{} + \ta{ரத்து} \textquotedblleft{}cancellation\textquotedblright{}.

\begin{cousinweb}
\textbf{\ta{விவாகரத்து}} (\emph{vivākarattu}) is a compound: \emph{vivāha}, \textquotedblleft{}marriage\textquotedblright{}, borrowed from Sanskrit (CDIAL 11920), and \ta{ரத்து}, \textquotedblleft{}cancellation\textquotedblright{}, which is usually traced to Arabic through Hindustani.
\end{cousinweb}

\begin{grammarlens}[title={What you've built}]
One more word for describing what you lived and felt.
\end{grammarlens}

\subsection*{Your turn}
\begin{itemize}
\raggedright
  \item \emph{Say it:} \emph{vivākarattu}
  \item \emph{Say it:} \emph{vivākarattu}, once more
  \item \emph{Say it:} \emph{pala varuṭac caṇṭaikkup piṟaku avarkaḷ vivākarattu ceytārkaḷ.} — After years of fighting they divorced.
  \item \emph{Recall:} say the Tamil for ill feeling, then the Tamil for to sulk, then say \emph{vivākarattu} again
\end{itemize}

\begin{tcolorbox}[breakable,colback=teal!4,colframe=teal!35!black,title={Before you move on}]
What does \ta{விவாகரத்து} mean? (\textquotedblleft{}A divorce\textquotedblright{}.) Say it once more.
\end{tcolorbox}
